% CRC Multi-Omics Drug-Target Atlas (main text)
% Target tier (rev1): Briefings in Bioinformatics / Journal of Translational Medicine
\documentclass[11pt,a4paper]{article}

\usepackage[utf8]{inputenc}
\usepackage[T1]{fontenc}
\usepackage[margin=2.5cm]{geometry}
\usepackage{graphicx}
\usepackage{booktabs}
\usepackage{longtable}
\usepackage{lineno}
\usepackage[hidelinks]{hyperref}
\usepackage[numbers,sort&compress]{natbib}
\usepackage{amsmath}
\usepackage{amssymb}
\usepackage{subcaption}
\usepackage{tikz}
\usetikzlibrary{arrows.meta,positioning,calc}
\captionsetup[subfigure]{labelformat=parens,labelsep=space}
\renewcommand{\thesubfigure}{\Alph{subfigure}}

\graphicspath{{results/figures_rev3/}{results/figures/}{results/phase6_scrna/}}

\title{Limited plasma proteomic coverage constrains genetic\\drug-target prioritization in colorectal cancer}

\author{Jiajie Zhou\textsuperscript{1,*}\\[0.4em]
\normalsize\textsuperscript{1}The Affiliated Huaian No.1 People's Hospital of\\ Nanjing Medical University,\\ Huaian, Jiangsu, China\\[0.3em]
\normalsize\textsuperscript{*}Corresponding author: Jiajie Zhou, MD.\\[0.15em] E-mail: zjjwcwk2025@njmu.edu.cn; Tel: +86-15805233415}

\begin{document}
\linenumbers

\maketitle

%==========================================================================
\section*{Abstract}
%==========================================================================

\noindent\textbf{Background:} Colorectal cancer (CRC) GWAS have identified over 100 susceptibility loci, but
the extent to which the resulting target lists can be validated at the protein
level has not been quantified.

\noindent\textbf{Methods:} We integrated blood cis-eQTL from eQTLGen
($n=31,684$), two plasma pQTL panels (UKB-PPP and deCODE, $n\approx35,000$ each),
a 78,473-case CRC GWAS, single-cell and spatial transcriptomics, and drug-safety
annotation into a single target-prioritization pipeline built on Mendelian randomization.

\noindent\textbf{Results:} From 15,582 gene-level
SMR tests, 75 probes surpassed Bonferroni significance ($p<3.21\times10^{-6}$)
and 43 of 76 probe--gene pairs (57\%) passed Bayesian colocalization (PPH4 $>$
0.8); 19 genes reached Tier 1 status through SuSiE fine-mapping, and six non-MHC
genes (UTP11, BMP2, PNKD, LAMC1, TMBIM1, GPBAR1) formed the core candidate set,
with BMP2 carrying the largest per-SD CRC risk effect (OR = 1.38, 95\% CI:
1.27--1.50) and GPBAR1 the only protective effect (OR = 0.78, 95\% CI:
0.73--0.83). At the protein layer, however, only 15 of the 62 protein-coding
genes had matching plasma protein measurements (17 in total, including two
TNF-pathway genes analyzed outside the Bonferroni set), only nine had cis-pQTL instruments
at genome-wide significance, and only one (CCM2, PPH4 = 0.9514) colocalized with
CRC risk, against 43 of 76 (57\%) at the eQTL layer; five of the six core
candidates were not measured in either panel. Instrument-strength and pleiotropy
diagnostics (all instruments $F>10$; MR-Egger intercepts non-significant for
CCM2 and LIMA1 only; MR-PRESSO outlier correction left the direction of every
estimate unchanged) indicate that the discordance is a property of the evidence
layers rather than of individual genes. 

\noindent\textbf{Conclusions:} Protein-layer coverage, not genetic
significance, is therefore the current limiting step in CRC drug-target
prioritization; we recommend reporting cis-pQTL coverage alongside target lists,
and we identify CCM2 and BMP2 as the candidates that remain defensible once this
constraint is stated explicitly.

\textbf{Keywords:} colorectal cancer, SMR, Mendelian randomization, colocalization,
pQTL, drug repurposing, spatial transcriptomics

%==========================================================================
\section{Introduction}
%==========================================================================

Colorectal cancer (CRC) accounts for approximately 1.9 million new cases and
935,000 deaths annually worldwide~\cite{sung2021global}, ranking third in incidence
and second in mortality among all cancers~\cite{bray2024global, siegel2026colorectal}.  Large-scale GWAS have
identified over 100 CRC susceptibility loci~\cite{huyghe2019discovery,
law2019association, fernandez2023deciphering}, yet the functional interpretation of
these loci and their translation into therapeutic targets remain significant
challenges~\cite{hazelwood2025multi, finan2017druggable, buniello2025opentargets}.

Summary-based Mendelian randomization (SMR) combined with Bayesian
colocalization has emerged as a powerful framework for prioritizing causal
genes at GWAS loci by integrating cis-eQTL data~\cite{zhu2016integration, wu2018integrative,
gusev2016integrative, giambartolomei2014, vosa2021eqtlgen, gtex2020}. However,
most efforts rely exclusively on transcript-level evidence, leaving the
protein-level validation---critical for drug target development---largely
unexplored.
The emergence of large-scale plasma proteomic GWAS datasets, including the UK
Biobank Pharma Proteomics Project (UKB-PPP) with Olink-based measurements and
deCODE genetics with SomaScan-based measurements, now enables systematic
pQTL-based drug target validation~\cite{sun2023plasma, ferkingstad2021large, pietzner2021mapping},
and Mendelian randomization of plasma proteins is now an established route to
causal target assessment~\cite{zheng2020phenome, schmidt2020drugtarget, hemani2018mrbase}.
Proteomic evidence is particularly valuable because most drug targets are
proteins, and transcript-level associations do not always propagate to the
protein level due to post-transcriptional and post-translational
regulation~\cite{vogel2012insights, liu2016impact}.
Whether these resources actually cover the genes that blood eQTL-based screens
nominate has not been quantified. Blood eQTL and plasma pQTL are measurements of
different molecules, and the degree to which the two agree at the same locus is
largely unknown.

We therefore set out not only to prioritize CRC drug targets from human genetic
data but to quantify how much of the resulting list can be tested at the protein
layer. The framework integrates six layers of evidence: (1) genome-wide SMR using blood cis-eQTL
from eQTLGen ($n=31,684$); (2) Bayesian colocalization and SuSiE
fine-mapping at the top loci; (3) dual pQTL validation using both UKB-PPP and
deCODE datasets; (4) independent directional consistency assessment in FinnGen R13 ($\sim$10,500 CRC
cases); (5) tissue-level characterization through single-cell RNA-seq and spatial
transcriptomics; and (6) drug annotation, repositioning scoring, and PheWAS
safety screening. Rather than reporting a ranked target list alone, we report the
coverage of each layer and the point at which the evidence stops.

%==========================================================================
\section{Methods}
%==========================================================================

\subsection{Study design overview}

This study employed a multi-phase design: (1) genome-wide SMR discovery using
blood cis-eQTL data; (2) SMR effect-size clinical interpretation through per-SD
odds ratio conversion; (3) Bayesian colocalization, SuSiE fine-mapping, and
independent locus definition; (4) dual pQTL Mendelian randomization (UKB-PPP and
deCODE); (5) pQTL colocalization; (6) directional consistency assessment in FinnGen;
(7) single-cell RNA-seq, spatial transcriptomics with microenvironment zoning
and gene-level cross-correlation; and (8) drug annotation with druggability
grading and PheWAS safety screening.

\subsection{GWAS summary statistics}

CRC GWAS summary statistics were obtained from GCST90255675 (78,473 cases,
107,143 controls, European ancestry), accessed through the NHGRI-EBI GWAS
Catalog~\cite{sollis2023gwas}. All genomic coordinates are reported
in hg38.

\subsection{eQTL data and SMR analysis}

Blood cis-eQTL data were obtained from eQTLGen ($n=31,684$)~\cite{vosa2021eqtlgen},
with tissue-level eQTL data from the Genotype-Tissue Expression (GTEx)
Project~\cite{gtex2020}. SMR analysis
was performed using the SMR software (v1.4.2)~\cite{zhu2016integration} with a cis window of $\pm$1 Mb
around each gene's transcription start site. The Bonferroni-corrected
significance threshold was $p < 3.21 \times 10^{-6}$ (0.05/15,582 tested
genes). HEIDI heterogeneity test was applied to distinguish linkage from
pleiotropy~\cite{wu2018integrative}. For clinical interpretation, SMR $b$ coefficients were converted
to odds ratios (OR = $\exp(b_{\text{SMR}})$) with 95\% confidence intervals
($\exp(b \pm 1.96 \times \text{se})$), representing the change in CRC risk per
one standard deviation increase in genetically predicted gene expression.

\subsection{Independent locus definition}

SMR-significant probes were collapsed into independent genetic loci using a
$\pm$1 Mb greedy clustering approach in genomic order. For each locus, genes
within 1 Mb of each other were grouped, and the locus span was defined as the
minimum and maximum gene positions within the cluster.

\subsection{Bayesian colocalization and SuSiE fine-mapping}

Colocalization analysis was performed using the coloc R package (v5.2.3)~\cite{giambartolomei2014},
with default priors ($p_1 = 10^{-4}$, $p_2 = 10^{-4}$, $p_{12} = 10^{-5}$).
The PPH4 threshold for colocalization was $>$ 0.8. Coloc was performed on
all 76 probe--gene pairs derived from the 75 Bonferroni-significant SMR
probes (one probe maps to two gene symbols). To assess sensitivity to
prior specification, we repeated coloc analyses varying $p_{12}$ across
$10^{-6}$, $10^{-5}$, and $10^{-4}$ for the six non-MHC Tier 1 genes. SuSiE
fine-mapping was applied using the susieR package (v0.14.2)~\cite{wang2020susie} with $L=10$ and
default parameters. A total of 27 coloc-passing genes were selected for SuSiE
analysis, comprising the six initial non-MHC coloc-passing candidates, all six
MHC-region genes, and fifteen additional non-MHC genes with high coloc PPH4
(covering the full spectrum of coloc-passing genes). Sensitivity analysis
was performed across 108 parameter combinations varying SNP inclusion
thresholds (200, 300, 500), $L$ values (5, 10, 20), and cis windows ($\pm$500
kb, $\pm$1,000 kb), of which 65 combinations (60\%) succeeded. The 43 failures
were concentrated in the chr2:219Mb region at large windows due to LD matrix
complexity exceeding 1000 Genomes resolution. Tier
classification: Tier 1 = coloc (PPH4 $>$ 0.8) + SuSiE (converged); Tier 2 =
coloc only; Tier 3 = SMR-significant only.

Genes that failed the HEIDI heterogeneity test ($p_{\text{HEIDI}} < 0.01$)
but passed coloc and SuSiE were retained in Tier 1, as colocalization evidence
under a single-causal-variant model provides independent support against
pleiotropy at the tested locus. HEIDI failure in the presence of coloc is
interpreted as evidence of heterogeneity not attributable to shared causal
variants, and such genes are flagged in Supplementary Table~S2 and Supplementary Fig.~S5. Accordingly,
HEIDI was treated as a pleiotropy flag, not a formal exclusion criterion
when coloc and SuSiE converged on a shared causal variant.

\subsection{pQTL data and Mendelian randomization}

UKB-PPP pQTL data (Olink, $n\approx35,000$) were accessed via Synapse
(project syn51364943)~\cite{sun2023plasma, sudlow2015ukbiobank}. deCODE pQTL data (SomaScan v4, $n\approx35,000$) were
obtained from the deCODE genetics proteomics download portal~\cite{ferkingstad2021large}. Cis-pQTL
instruments were selected at $p < 5\times10^{-8}$ within $\pm$1 Mb of the
transcription start site. Two effect estimates are reported in this paper and
are not interchangeable. The single-variant Wald ratio uses the lead cis-pQTL
variant only; it is the estimate given in Table~\ref{tab:pqtl_coverage}(B) and the
one used for the protein-layer classification in
Fig.~\ref{fig:eqtl_pqtl_class}. The number of cis-pQTL instruments listed for a
gene in Table~\ref{tab:pqtl_coverage}(A) is the number of variants available at that
locus, not the number that contributes to the Wald ratio. The IVW estimate uses
all available cis-pQTL instruments and, with its sensitivity analyses, is
reported in Table~\ref{tab:mr_sensitivity}. Wherever the multi-instrument estimate
could be computed we report both, because the two can differ in magnitude and,
at one locus (CDKN1A), in sign. Mendelian randomization was performed using
Wald ratio (single instrument) and inverse-variance weighted (IVW; multiple
instruments) methods~\cite{burgess2013ivw}, with the weighted median~\cite{bowden2016median},
weighted mode~\cite{hartwig2017mode}, MR-Egger~\cite{bowden2015egger} and
MR-PRESSO~\cite{verbanck2018presso} estimators used as sensitivity analyses and
TwoSampleMR~\cite{hemani2018mrbase} used for implementation. Reporting follows
the STROBE-MR guidance~\cite{skrivankova2021strobe}. For colocalization with quantitative pQTL data, minor
allele frequency from the pQTL dataset was supplied to estimate the trait
variance (sdY) as required by coloc v5.2.3.

The three core instrumental-variable assumptions were addressed as follows:
relevance was met by restricting instruments to cis variants associated with the
exposure at $p < 5\times10^{-8}$ in the source pQTL study; independence and the
exclusion restriction cannot be tested directly, so instruments were restricted
to cis regions to limit the scope for trans-acting pleiotropy, and we applied
HEIDI for the eQTL layer and MR-Egger, weighted median, weighted mode and
MR-PRESSO for the pQTL layer. This study was a secondary analysis of publicly
available summary statistics and was not pre-registered. Variants present in the exposure but not in the outcome dataset (or vice versa) were excluded rather than imputed. Because the UKB-PPP panel and the CRC GWAS both include UK Biobank participants, the exposure and outcome samples may partially overlap, which can bias two-sample MR estimates towards the confounded observational estimate.

Directionality was assessed with the Steiger test~\cite{hemani2017steiger}. For
every cis-pQTL instrument we compared the variance in plasma protein level
explained by the variant ($r^2$ derived from the pQTL effect size and standard
error at the panel sample size; $n = 35{,}278$ for deCODE and $n = 33{,}533$ for
UKB-PPP) with the variance in CRC risk explained by the same variant
($N = 185{,}616$), and we additionally ran the per-gene aggregate test
implemented in \texttt{TwoSampleMR::directionality\_test}. Because the aggregate
statistic sums $r^2$ across correlated variants and can exceed unity, a second
per-gene test restricted to the lead cis-pQTL variant was also performed
(Supplementary Table~S10).


\subsection{FinnGen directional consistency assessment}

Directional consistency was assessed in FinnGen R13~\cite{kurki2023finngen} C3\_COLORECTAL (10,500
cases) by comparing discovery SMR effect directions with FinnGen MR estimates;
this cohort was treated as a power-limited consistency check, not a
formal replication sample.

\subsection{Single-cell and spatial transcriptomics}

CRC single-cell RNA-seq data were analyzed using Seurat~\cite{hao2021seurat}. Cell-type
specificity was assessed by comparing expression across annotated cell types.
Spatial transcriptomics data (Visium)~\cite{stahl2016spatial} were processed using Seurat spatial
workflow, and the fraction of spots with detectable expression was recorded for
each candidate gene. To characterize the tumor
microenvironment, Visium spots were partitioned using AddModuleScore
cell-type signature scoring followed by K-means clustering. The optimal number
of spatial microenvironment zones was selected by the elbow method on
within-cluster sum of squares, yielding $k=6$ zones. Spatial gene-level
cross-correlation was assessed by computing pairwise Spearman correlations
between Tier 1 or coloc-passing gene expression values across Visium spots
within the same chromosome locus, with a cross-correlogram function evaluating
correlation decay across spot-to-spot distances in 30 bins of 100 $\mu$m.

\subsection{Drug annotation and PheWAS}

Drug-target associations, target tractability and safety annotations were
retrieved from the Open Targets Platform~\cite{buniello2025opentargets} and
Pharos~\cite{kelleher2023pharos}, with druggable-genome classification
following Finan et al.~\cite{finan2017druggable} and phenome-wide association
results interpreted following Denny et al.~\cite{denny2010phewas}.

Drug annotations were obtained from DrugBank (v5.1) and Open Targets Platform
(v24.06). Repositioning scores were calculated as a weighted sum of genetic
evidence (SMR + coloc + SuSiE + HEIDI), biological evidence (TCGA expression
+ cell-type specificity), and drug availability. Target tractability was
assessed using the Illuminating the Druggable Genome (IDG) classification
(Tclin/Tchem/Tbio/Tdark), structural data availability (PDB), and chemical
probe availability (ChEMBL/DrugBank). PheWAS safety screening used Open
Targets Genetics GWAS associations to identify pleiotropic disease
associations for each Tier 1 target gene; these associations reflect general
population genetic correlations and should be interpreted as hypothesis-generating
safety signals, not regulatory-grade pharmacovigilance evidence.

\subsection{Competitor comparison}

Our set of 69 unique gene symbols (from 75 unique genes: 62 protein-coding
and 13 non-coding) was systematically
compared against six prior CRC genetic studies: four low-tier SMR/MR studies
(Hong et al.~2025, 9 genes~\cite{hong2025integrative}; Wang et al.~2025, 4
genes~\cite{wang2025discovoncol}; Xia and Wang, BMC Cancer, 6
genes~\cite{xia2025bmccancer}; Tian et al.~2025, 8 genes~\cite{tian2025dbi}) and two high-tier Nature Communications studies (Hazelwood
2025, 41 genes from multi-tissue TWAS~\cite{hazelwood2025multi}; Chen et
al.~2024, 290 published genes from 254K-individual GWAS + SMR + multi-tissue
eQTL/mQTL~\cite{chen2024integrative}, the largest prior CRC fine-mapping study).
Chen 2024 gene lists were extracted from Supplementary Data 20 (288 putative
target genes from functional genomics + coloc), Supplementary Data 13 (154
cis-eQTL-significant genes), and Supplementary Data 11--12 (48 + 22 in-silico
predicted genes), merged and deduplicated to 290 entries, of which 289 could be resolved to a unique
gene symbol; this 289-symbol set was used for all overlap calculations. Overlap
was defined as identical or synonymous gene symbols. At the locus level, a
locus was considered overlapping with Chen 2024 if it contained at least one gene
from the Chen 2024 gene set. We note that SMR (blood eQTL) and TWAS (multi-tissue
GTEx + sQTL) are methodologically distinct and their gene sets are not expected
to fully overlap, a consideration acknowledged in the interpretation of overlap
statistics.

%==========================================================================


\section{Results}
%==========================================================================

\subsection{A resolved genetic map of CRC susceptibility genes}

The overall study design and the two evidence layers are summarized in
Fig.~\ref{fig:design}, encompassing seven analytical phases from SMR discovery
through drug target prioritization.

\begin{figure}[htbp]
\centering
\begin{subfigure}{\textwidth}
\centering
\includegraphics[width=\linewidth]{Fig1a_conceptual_overview.pdf}
% The diagram above is exported to Fig1a_conceptual_overview.pdf by
% scripts/73_export_fig1a_panel.py. The TikZ source it was built from is kept
% here, inactive, for future edits.
% \begin{tikzpicture}[
%   font=\sffamily,
%   box/.style={draw=black!60, rounded corners=1.5pt, align=center, inner sep=3pt,
%               text width=31mm, minimum height=8mm, font=\footnotesize},
%   head/.style={align=center, font=\footnotesize\bfseries, text width=40mm},
%   arw/.style={-{Latex[length=1.6mm,width=1.4mm]}, draw=black!70, line width=0.5pt}
% ]
% \node[head] (hA) at (0,0)    {Genetic layer\\[1pt]{\mdseries\scriptsize blood cis-eQTL}};
% \node[head] (hB) at (5.6,0)  {Protein layer\\[1pt]{\mdseries\scriptsize plasma pQTL}};
% \node[head] (hC) at (11.2,0) {Interpretation};
%
% \node[box, fill=blue!7]     (a1) at (0,-1.1) {15{,}582 genes\\tested};
% \node[box, fill=blue!7]     (a2) at (0,-2.6) {75 probes significant\\$p<3.21\times10^{-6}$};
% \node[box, fill=blue!7]     (a3) at (0,-4.1) {\textbf{43/76 (57\%)}\\colocalised};
% \node[box, fill=blue!7]     (a4) at (0,-5.6) {19 genes Tier 1\\(SuSiE-converged)};
% \draw[arw] (a1)--(a2); \draw[arw] (a2)--(a3); \draw[arw] (a3)--(a4);
%
% \node[box, fill=orange!12]  (b1) at (5.6,-1.1) {62 protein-coding\\candidates};
% \node[box, fill=orange!12]  (b2) at (5.6,-2.6) {15 measured in plasma\\(17 with TNF pathway)};
% \node[box, fill=orange!12]  (b3) at (5.6,-4.1) {9 with cis-pQTL\\$p<5\times10^{-8}$};
% \node[box, fill=orange!12]  (b4) at (5.6,-5.6) {\textbf{1 colocalised}\\CCM2 (PPH4 = 0.95)};
% \draw[arw] (b1)--(b2); \draw[arw] (b2)--(b3); \draw[arw] (b3)--(b4);
%
% \draw[arw] (a1.east) -- (b1.west)
%   node[midway, above, font=\scriptsize, align=center] {protein-coding};
% \draw[arw, dashed] (a3.east) -- (b4.west)
%   node[midway, below, font=\scriptsize\itshape, align=center] {coverage collapse};
%
% \node[box, fill=gray!8, text width=40mm, minimum height=30mm] (c1) at (11.2,-3.35)
%   {\textbf{Limiting step}\\[3pt] protein-layer coverage,\\not genetic significance\\[5pt]
%    eQTL $\neq$ pQTL\\[3pt] 5 of 6 core candidates\\not measured in either panel};
% \draw[arw] (b3.east) -- (c1.west);
% \end{tikzpicture}
\caption{Conceptual overview of the two evidence layers. The genetic layer (left) is
saturated: 15{,}582 gene-level SMR tests yield 43 colocalizing probe--gene pairs out of 76.
The protein layer (center) collapses: only 15 of the 62 protein-coding candidates have plasma
measurements, nine have cis-pQTL instruments, and one colocalises with CRC risk. The limiting
step is therefore protein-layer coverage, not genetic significance, and blood eQTL and plasma
pQTL evidence are not interchangeable.}
\label{fig:design_concept}
\end{subfigure}

\begin{subfigure}{\textwidth}
\centering
\includegraphics[width=\linewidth]{Fig1b_study_design.pdf}
\caption{Study design and analytical pipeline. Multi-omics drug-target discovery workflow
integrating eQTL SMR discovery, colocalization and SuSiE fine-mapping, pQTL Mendelian
randomization, FinnGen replication, GTEx Colon tissue sensitivity, single-cell and Visium
spatial validation, and drug annotation. Numbers refer to the results reported here:
75 genes passing Bonferroni-corrected SMR (62 protein-coding), 43 of 76 colocalizing
probe--gene pairs, one convergent gene at the protein layer, one of eight Tier 1 genes
replicating in FinnGen, eight genes with known pharmacological agents, and six core Tier 1
genes characterized in detail.}
\label{fig:design_panel}
\end{subfigure}

\vspace{0.5em}

\caption{Study design, evidence layers and genetic-layer colocalization screen.}
\label{fig:design}
\end{figure}

\begin{figure}[tbp]
\ContinuedFloat
\centering
\setcounter{subfigure}{2}

\begin{subfigure}{\textwidth}
\centering
\includegraphics[width=\linewidth]{Fig1c_coloc_pph4.pdf}
\caption{Bayesian colocalization of the 76 probe--gene pairs.
Posterior probability of colocalization (PPH4, coloc.abf) shown for the 49 pairs with PPH4 $>$ 0.5:
43 pairs (57\%) passed the colocalization threshold (PPH4 $>$ 0.8, dark blue) and 6 were marginal
(0.5--0.8, light blue). The remaining 27 pairs (6 weak, 21 failing) are omitted from the panel.}
\label{fig:coloc}
\end{subfigure}
\end{figure}

We performed SMR using cis-eQTL data from eQTLGen ($n=31,684$ blood
samples) and CRC GWAS summary statistics from GCST90255675 (78,473 cases,
107,143 controls). Of 15,582 genes tested with at least one cis-eQTL instrument
(at $p < 5\times10^{-8}$ within $\pm$1 Mb of the transcription start site), 75
passed the Bonferroni-corrected significance threshold ($p < 3.21 \times
10^{-6}$; 0.05/15,582). The genomic inflation factor was $\lambda = 1.600$, consistent with LD-coupled cis-eQTL SMR statistics~\cite{wainberg2019} (Supplementary Fig.~S1).

The top 10 SMR associations are listed in Supplementary Table~S1. The
top signal was \textit{CASC8} (cancer susceptibility 8; $p =
3.36\times10^{-28}$, $b_{\text{SMR}} = 0.236$), a lncRNA within the
8q24 CRC risk locus. Among protein-coding genes, the top hits included
\textit{LAMC1} ($p = 7.14\times10^{-24}$), \textit{LIMA1} ($p =
1.44\times10^{-19}$), \textit{BMP2} ($p = 5.77\times10^{-15}$), and
\textit{TMBIM1} ($p = 5.03\times10^{-14}$).

GENCODE v47 annotation of the 75 unique genes identified 62 protein-coding
genes and 13 non-coding genes (8 lncRNAs, 1 snRNA, 2 pseudogenes, and 2
transcripts absent from GENCODE v47). The 62 protein-coding genes
formed the primary candidate set for downstream colocalization and
validation (Supplementary Table~S2). Seven non-coding genes lacked a gene
symbol, and one probe mapped to two distinct gene symbols. The 75 genes
collapsed to 69 unique gene symbols, and 76 probe--gene pairs moved
forward to colocalization analysis.

Across the 75 SMR-significant genes we tested colocalization on 76 probe--gene pairs
(\textit{UTP11} contributed two probes at chr1:38.4~Mb, hence 76 pairs for 75 unique
genes; Fig.~\ref{fig:coloc}). 43 of 76 pairs (57\%) showed strong evidence of shared causal variants
(PPH4 $>$ 0.8). The four strongest independent loci were \textit{SMAD7} (PPH4 = 0.9997),
\textit{UTP11} (PPH4 = 0.9996; two concordant probes), \textit{TNF} (PPH4 = 0.9991)
and \textit{BMP2} (PPH4 = 0.9971), closely followed by \textit{B9D2} (PPH4 = 0.9967).
\textit{SMAD7}, the highest-PPH4 non-MHC hit, was
not resolved by SuSiE because the 1000 Genomes European reference LD matrix did
not converge at chr18:46,461,652; it is therefore a high-confidence
colocalization-only candidate, not part of the Tier 1 set defined below. Six of
the 43 colocalization-passing genes reside
within the MHC region (chr6:29--33 Mb; \textit{TNF}, \textit{LTA}, \textit{HLA-K},
\textit{HLA-DRB5}, \textit{ZFP57}, \textit{HCG27}), where extended LD could inflate
coloc PPH4. Applying LD-aware SuSiE fine-mapping to this region returned
single-SNP credible sets (PIP = 1.0) for all six genes. However, several of
these credible sets were driven by overlapping variants---notably \textit{TNF},
\textit{LTA}, and \textit{HCG27} were resolved to essentially the same lead SNPs---so
the MHC signals do not represent independent causal variants and are interpreted
conservatively here.

Conversely, 21 probe--gene pairs showed essentially no evidence of colocalization
(PPH4 $<$ 0.01), including \textit{CASC8} (PPH4 = $9.2\times10^{-60}$),
\textit{UTP23} (PPH4 = $1.2\times10^{-28}$), \textit{LIMA1} (PPH4 =
$2.7\times10^{-7}$), and \textit{COX14} (PPH4 = $6.5\times10^{-12}$), primarily
due to high PPH3 indicating independent causal variants at the same locus.
Six genes had marginal colocalization (PPH4 0.5--0.8) and six showed weak
evidence (PPH4 0.01--0.5).

To resolve potential multiple causal signals and to sanity-check coloc against LD,
we applied SuSiE fine-mapping (susie\_rss, 1000 Genomes Phase 3 European LD
reference, $L$ up to 10 for non-MHC and up to 20 for MHC loci; coverage = 0.9) to
27 protein-coding coloc-passing loci: the six prioritized non-MHC genes (UTP11,
BMP2, PNKD, LAMC1, TMBIM1, GPBAR1), all six MHC-region genes (\textit{TNF},
\textit{LTA}, \textit{HLA-K}, \textit{HLA-DRB5}, \textit{ZFP57}, \textit{HCG27}),
and the fifteen highest-PPH4 remaining non-MHC protein-coding loci. Five
non-coding (Ensembl-only) probes and eleven lower-ranked protein-coding loci among
the 43 coloc-passing pairs were not fine-mapped. Of these 27 loci, 19 converged to
a credible set (six MHC and thirteen non-MHC genes), yielding 19 Tier 1 genes in
total, each with single-SNP credible sets (PIP = 1.0; Supplementary Table~S3).
Eight of the fifteen non-MHC loci (SMAD7, B9D2, MMP11, TMEM258, SMAD6, CCDC97,
STAT6, and PGAP3) were not resolved because the reference LD matrix failed to
converge at these loci, a technical limitation, not a biological one.
The number of reported credible sets scaled with the user-specified $L$ (e.g., UTP11
returned 5, 10 and 15 credible sets as $L$ was raised), and every resolved locus
returned a single-SNP credible set with PIP = 1.0---both expected when a reference
LD matrix is too coarse to separate the causal signal. The single-SNP credible sets
are therefore consistent with, but do not by themselves establish, a single causal
variant at each locus; we accordingly interpret SuSiE as an LD-aware consistency
check on colocalization rather than independent proof of a single causal variant.
All 19 genes were classified as Tier 1, defined as coloc PPH4 $>$ 0.8 with a
converged SuSiE credible set; SuSiE convergence is treated here as an LD-aware
consistency criterion rather than causal resolution. The complete Tier 1 set
comprises 13 non-MHC genes (UTP11, BMP2, PNKD, LAMC1, TMBIM1, GPBAR1, CCM2,
SF3A3, CLDN4, LEMD3, METTL27, GDPGP1, ABHD12B) and 6 MHC-region genes (TNF, LTA,
HLA-K, HLA-DRB5, ZFP57, HCG27); the six MHC genes are counted nominally, because
their credible sets were driven by overlapping variants and therefore do not
represent independent causal signals (see above). Among these, the six non-MHC
Tier 1 genes were selected for in-depth downstream characterization and are
summarized in Table~\ref{tab:tier1} (UTP11, BMP2, PNKD, LAMC1, TMBIM1, GPBAR1). Regional association plots for all Tier 1 loci are provided in Supplementary Fig.~S2. Among the
remaining 13 Tier 1 genes, CCM2 merits particular attention as the only Tier 1
gene with convergent pQTL colocalization evidence---characterized in detail in the
Discussion (subsection \emph{eQTL--pQTL colocalization discordance})---while the
remaining 12 are described in the Supplementary Materials. Six of the 19 Tier 1
genes (CCM2, GPBAR1, LAMC1, LTA, PNKD, TMBIM1) were previously nominated as CRC
drug-target candidates in the comparator studies we surveyed, whereas the
remaining 13---including eight non-MHC protein-coding genes---were not prioritized as candidate targets by those studies; BMP2 appears in the broader Chen 2024 gene set but not among their prioritized targets.

\begin{table}[htbp]
\centering
\caption{The six core non-MHC Tier 1 genes with colocalization and fine-mapping support.}
\label{tab:tier1}
\small
\resizebox{\textwidth}{!}{%
\begin{tabular}{lllllllll}
\toprule
  Gene & Chr & PPH4 & $n_{\text{CS}}$ & PIP & $p_{\text{SMR}}$ & Known drug & Risk categories & High-risk phenotypes \\
\midrule
  UTP11  & 1 & 0.9996 & 10 & 1.00 & $1.56\times10^{-13}$ & --- & 4 & 6 \\
  BMP2   & 20 & 0.9971 & 10 & 1.00 & $5.77\times10^{-15}$ & --- & 2 & 2 \\
  PNKD   & 2 & 0.9337 & 2 & 1.00 & $2.38\times10^{-15}$ & --- & 4 & 8 \\
  LAMC1  & 1 & 0.9294 & 10 & 1.00 & $7.14\times10^{-24}$ & OCRIPLASMIN (approval) & 3 & 3 \\
  TMBIM1 & 2 & 0.8262 & 3 & 1.00 & $5.03\times10^{-14}$ & --- & 3 & 7 \\
  GPBAR1 & 2 & 0.8132 & 3 & 1.00 & $1.75\times10^{-14}$ & --- & 2 & 8 \\ 
\bottomrule
\end{tabular}
}

\tblnote{Colocalization is defined as PPH4 $>$ 0.8 and fine-mapping support as convergence of SuSiE to a credible set. Risk categories and high-risk phenotypes are from Phenome-Wide Association Study (PheWAS) screening in the Open Targets Platform (see Methods and Supplementary Table~S8): risk categories = number of distinct organ-system categories carrying at least one genetic association for the gene; high-risk phenotypes = number of non-CRC phenotypes with genetic evidence in predefined high-risk categories.}
\end{table}


The chr2:219Mb region, containing \textit{PNKD}, \textit{TMBIM1}, and
\textit{GPBAR1}, was resolved into eight single-SNP credible sets
(PNKD: 2, TMBIM1: 3, GPBAR1: 3). To test whether this singleton pattern reflected
the SNP-pruning and $L$ settings rather than the underlying genetic architecture,
we performed a sensitivity sweep across 108 parameter combinations (SNP caps of
200/300/500, $L$ of 5/10/20, and $\pm$500 kb/$\pm$1,000 kb cis windows) over the six
prioritized genes. Of these, 65 combinations (60\%) succeeded and 43 failed, the
latter concentrated in the chr2:219Mb region at $\pm$1,000 kb windows, where LD
complexity exceeded the resolution of the 1000 Genomes Phase 3 European reference
panel. All successful combinations returned single-SNP credible sets with
PIP = 1.0, but---as noted above---the number of credible sets tracked the
user-specified $L$, and we therefore regard the singleton pattern as a property of
the method/settings, not evidence establishing a single causal variant (Supplementary Fig.~S3).

We used the coloc.abf default priors ($p_1 = p_2 = 10^{-4}$, $p_{12} = 10^{-5}$);
a prior sensitivity analysis showed that PPH4 values for the six non-MHC Tier 1
genes remained stable across $p_{12}$ from $10^{-6}$ to $10^{-4}$ (Supplementary
Table~S3).

SMR $b$ coefficients convert to odds ratios (OR = $\exp(b_{\text{SMR}})$),
the change in CRC risk per SD increase in genetically predicted gene
expression (Supplementary Table~S1 lists $b_{\text{SMR}}$ and $p_{\text{HEIDI}}$
for the ten strongest associations). Six non-MHC
genes were subsequently prioritized by colocalization and SuSiE. Five increased CRC risk per SD:
BMP2 (OR = 1.38, 95\% CI: 1.27--1.50), TMBIM1
(OR = 1.35, 95\% CI: 1.25--1.46), UTP11
(OR = 1.27, 95\% CI: 1.20--1.36), LAMC1 (OR = 1.14, 95\% CI:
1.11--1.17), and PNKD (OR = 1.12, 95\% CI: 1.09--1.15). GPBAR1
was the only protective gene (OR = 0.78, 95\% CI: 0.73--0.83; per-SD odds
ratios for the six genes are shown in Supplementary Fig.~S4).

Six of the 18 Tier 1 genes with HEIDI results failed the HEIDI heterogeneity
test ($p_{\text{HEIDI}} < 0.01$, the standard SMR heterogeneity threshold) and
are flagged in Supplementary Table~S2 and Supplementary Fig.~S5. BMP2 had no HEIDI result; the remaining 12
passed.

To contextualize the genetic architecture, we collapsed the 75 SMR-significant
genes into genomic loci by greedy $\pm$1 Mb physical-distance clustering ordered
by top-SNP position (a distance-based definition, not an LD-$r^{2}$ threshold).
This identified 37 loci, of which 20 (54\%) contained a single gene and 17 (46\%)
harbored two or more genes (max: 8 genes at chr2:219Mb); the count was stable
under $\pm$500 kb and $\pm$250 kb windows. The chr2:219Mb cluster (222 kb, 8 genes
total; 5 coloc-passing: PNKD, TMBIM1, GPBAR1, AAMP, and one non-coding probe; 3
coloc-failing: ARPC2, CXCR1, CXCR2) and chr12:50.4Mb cluster (799 kb, 7 genes)
represented the most complex architectures.

\subsection{The protein layer does not corroborate the genetic layer}

The protein layer, however, covers only a small fraction of the genes that the
genetic layer prioritizes. Of the 62 protein-coding SMR-significant genes, 15 had
a matching protein measurement in at least one of the two plasma proteomic panels
(5 in UKB-PPP, 14 in deCODE, with 4 genes measured in both). Two further genes,
TNFRSF1A and TNFRSF1B, were measured in deCODE and carried into the protein-layer
analysis as members of the TNF pathway even though they did not reach the
Bonferroni threshold, giving a protein-layer analysis set of 17 genes
(Fig.~\ref{fig:coverage_funnel}). Of these 17 measured genes, only nine had cis-pQTL
instruments at genome-wide significance ($p < 5\times10^{-8}$) and could therefore
be tested by Mendelian randomization: four in UKB-PPP and five in deCODE (seven of
them within the 62-gene Bonferroni set). The remaining 47 protein-coding
genes---including five of the six core Tier 1 candidates---were not measured in
either panel and could not be tested at the protein layer at all.

To evaluate protein-level evidence, we performed pQTL Mendelian randomization
using two independent proteomic datasets. Coverage, instrument strength, and
colocalization results for all 17 measured genes are summarized in
Table~\ref{tab:pqtl_coverage}. From UKB-PPP, 5 of the 62
protein-coding genes had matching Olink protein measurements (CDKN1A, LTA,
NCF2, TET2, TNF). By the single-variant Wald ratio, MR analysis identified
CDKN1A (Wald $p = 0.020$) and TNF (Wald $p = 5.0\times10^{-4}$) as significant,
with effects in the risk-increasing direction for CDKN1A (Wald $b = +0.130$) and
the protective direction for TNF (Wald $b = -0.295$). The multi-instrument IVW
estimates agreed for TNF ($b = -0.062$, $p = 8.1\times10^{-29}$) but not for
CDKN1A, whose sign reversed ($b = -0.131$, $p = 1.3\times10^{-18}$). LTA and NCF2
were non-significant by the Wald ratio ($p = 0.506$ and $0.307$); by IVW the LTA
estimate became significant ($b = -0.022$, $p = 5.4\times10^{-44}$) while NCF2
remained null ($b = -0.033$, $p = 0.156$; Table~\ref{tab:mr_sensitivity}). However, none of the 5 genes showed pQTL--GWAS colocalization (max PPH4 = 0.275 for TNF; Supplementary Table~S4; Supplementary Fig.~S6).

From deCODE genetics, 16 unique genes (14 of them within the 62-gene Bonferroni set) had SomaScan protein measurements,
representing 12 deCODE-only genes beyond the 4 genes also covered by UKB-PPP
(CDKN1A, NCF2, TNF, LTA; analyzed above). Of the 12 deCODE-only genes, 5
carried cis-pQTL instruments at genome-wide significance ($p < 5\times10^{-8}$):
LIMA1 (70 instruments), CCM2 (94), STAT6 (61), TNFRSF1A (25) and TNFRSF1B (73).
By the single-variant Wald ratio, three were significant---LIMA1
($b = +0.858$, $p = 4.29\times10^{-7}$), CCM2 ($b = -0.766$, $p = 6.28\times10^{-6}$)
and STAT6 ($b = -0.029$, $p = 3.86\times10^{-5}$)---whereas TNFRSF1A
($b = -0.070$, $p = 0.349$) and TNFRSF1B ($b = -0.035$, $p = 0.614$) were not.
The multi-instrument IVW estimates agreed in magnitude for the two strongest
signals, CCM2 ($b = -0.515$, $p = 2.7\times10^{-17}$) and LIMA1
($b = +0.882$, $p = 1.3\times10^{-26}$), but not for the other three: STAT6 was
not significant by IVW ($b = -0.007$, $p = 0.427$) whereas TNFRSF1B was
($b = -0.052$, $p = 3.0\times10^{-7}$), and TNFRSF1A remained null by both
estimators ($b = -0.013$, $p = 0.637$). The seven remaining deCODE-only genes had
no cis-pQTL instruments at genome-wide significance, reflecting the
smaller sample size and higher technical variability of plasma proteomic
measurements.

Among the 5 deCODE genes with instruments, only CCM2 showed
evidence of pQTL--GWAS colocalization (PPH4 = 0.9514, 6,850 SNPs;
Supplementary Table~S5, Supplementary Fig.~S7). This makes CCM2 the only gene in our study with
convergent eQTL and pQTL colocalization evidence. CCM2 is also one
of the 19 Tier 1 genes (eQTL coloc PPH4 = 0.9633, SuSiE converged with
single-SNP credible set), making it the sole gene with triple-convergent
evidence across eQTL fine-mapping, pQTL MR, and pQTL colocalization. The
remaining 4 genes did not colocalize. STAT6, LIMA1, and TNFRSF1A showed
PPH3 $\approx 1.0$, indicating distinct causal variants for the pQTL and
GWAS signals; STAT6 and LIMA1 had significant pQTL MR associations (Wald
$p = 3.86\times10^{-5}$ and $4.29\times10^{-7}$), whereas TNFRSF1A did not
(Wald $p = 0.349$). TNFRSF1B showed no colocalization (PPH4 = 0.0084) and no
distinct pQTL signal (PPH3 = 0.287), yet its multi-instrument IVW estimate was
significant ($b = -0.052$, $p = 3.0\times10^{-7}$) even though its single-variant
Wald ratio was not ($p = 0.614$)---the opposite pattern to CDKN1A, and a further
reminder that the two estimators can disagree in either direction.

\begin{table}[htbp]
\centering
\caption{Plasma protein coverage and pQTL--GWAS colocalization for the 17 genes with matching protein measurements. Cis-pQTL instruments were defined at $p < 5\times10^{-8}$ within $\pm$1 Mb of the gene; $n$ is the number of instruments used in the Mendelian randomization analysis. Wald-ratio estimates are the ratio of the CRC GWAS effect to the cis-pQTL effect on protein level. PPH4, posterior probability of a shared causal variant between the cis-pQTL and CRC GWAS signals (coloc.abf); PPH4 $>$ 0.8 indicates colocalization. Median $F$ is the median instrument $F$ statistic across the cis-pQTL instruments. No genome-wide significant cis-pQTL was detected for the seven deCODE-only proteins or for TET2; TET2 was nonetheless tested by colocalization and showed no evidence of a shared causal variant (PPH4 $= 0.006$). Panel names refer to the proteomic panel in which the protein was measured; MR panel is the panel supplying the cis-pQTL instruments.}
\label{tab:pqtl_coverage}
\scriptsize
\resizebox{\textwidth}{!}{%
\begin{tabular}{llllllll}
\toprule
  Gene & Panel & $n$ cis-pQTL & Median $F$ & Wald $b$ ($p$) & PPH4 & Protein-level support \\
\midrule
  CCM2      & deCODE         & deCODE   & 94     & 30.4   & -0.766 ($6.28\times10^{-6}$) & 0.9514     & colocalizes \\
  TNF       & both           & UKB-PPP  & 2288   & 49.2   & -0.295 ($5.22\times10^{-4}$) & 0.2751     & no coloc \\
  LTA       & both           & UKB-PPP  & 10213  & 175.4  & -0.014 (0.506)       & 0.2605     & no coloc \\
  TNFRSF1B  & deCODE         & deCODE   & 73     & 127.5  & -0.035 (0.614)       & 0.0084     & no coloc \\
  LIMA1     & deCODE         & deCODE   & 70     & 32.8   & +0.858 ($4.29\times10^{-7}$) & $1.93\times10^{-4}$ & no coloc \\
  NCF2      & both           & UKB-PPP  & 28     & 62.6   & -0.071 (0.307)       & $2.14\times10^{-6}$ & no coloc \\
  STAT6     & deCODE         & deCODE   & 61     & 66.0   & -0.029 ($3.86\times10^{-5}$) & $4.65\times10^{-8}$ & no coloc \\
  CDKN1A    & both           & UKB-PPP  & 228    & 66.9   & +0.130 (0.020)       & $1.33\times10^{-9}$ & no coloc \\
  TNFRSF1A  & deCODE         & deCODE   & 25     & 54.2   & -0.070 (0.349)       & $1.58\times10^{-10}$ & no coloc \\
  ARPC5     & deCODE         & ---      & none   & ---    & ---                  & ---        & no cis-pQTL instruments \\
  BMP2      & deCODE         & ---      & none   & ---    & ---                  & ---        & no cis-pQTL instruments \\
  CABLES2   & deCODE         & ---      & none   & ---    & ---                  & ---        & no cis-pQTL instruments \\
  CERS5     & deCODE         & ---      & none   & ---    & ---                  & ---        & no cis-pQTL instruments \\
  MZF1      & deCODE         & ---      & none   & ---    & ---                  & ---        & no cis-pQTL instruments \\
  STAB1     & deCODE         & ---      & none   & ---    & ---                  & ---        & no cis-pQTL instruments \\
  UBE2M     & deCODE         & ---      & none   & ---    & ---                  & ---        & no cis-pQTL instruments \\
  TET2      & UKB-PPP        & ---      & none   & ---    & ---                  & 0.0059     & no cis-pQTL instruments \\ 
\bottomrule
\end{tabular}
}
\end{table}


Formal instrument-strength and pleiotropy diagnostics for all nine genes with
cis-pQTL instruments are reported in Table~\ref{tab:mr_sensitivity} and
Supplementary Figs.~S8--S13. Every instrument across the nine genes had an F
statistic well above the conventional weak-instrument threshold (minimum
F = 29.7; median F per gene 30.4--175.4). Cochran's Q indicated substantial
heterogeneity in the eight genes with more than two instruments (TNFRSF1A
Q = 94.8, df = 24, p = 2.3e-10; NCF2 Q = 103.1, df = 27, p = 7.9e-11; STAT6
Q = 177.0, df = 60, p = 1.7e-13; LIMA1 Q = 575.7, df = 69, p = 5.3e-81; CCM2
Q = 770.7, df = 93, p = 2.8e-107; CDKN1A Q = 687.1, df = 227, p = 8.8e-48; TNF
Q = 9683.2, df = 2287, p < 1e-300; LTA Q = 63797, df = 10212, p < 1e-300) but
not in TNFRSF1B, whose 73 instruments were homogeneous (Q = 43.1, df = 72,
p = 0.997). MR-Egger intercepts were non-significant for CCM2 (b = 0.013,
p = 0.55) and LIMA1 (b = 0.074, p = 0.11), the two genes with colocalizing or
near-colocalizing protein-level support, and significant for the remaining seven
genes (TNFRSF1A b = -0.017, p = 5.1e-4; NCF2 b = +0.041, p = 0.034; STAT6
b = 0.021, p = 3.0e-12; TNFRSF1B b = -0.014, p = 8.7e-5; CDKN1A b = -0.012,
p = 1.8e-4; TNF b = -0.012, p < 1e-300; LTA b = -0.002, p = 7.4e-4). MR-PRESSO
was run at NbDistribution = 5000 for the seven genes with at most 228
instruments and returned a significant global test for every one of them
(p < 2.0e-4, the resolution limit of the 5,000-draw null distribution); the
number of instruments flagged as outliers and the resulting outlier-corrected
estimates are given in Table~\ref{tab:mr_sensitivity}(C). Removal of the
flagged instruments left the two strongest protein-level signals intact and
larger in magnitude (CCM2, b = -0.52 raw versus -0.76 outlier-corrected,
p = 1.9e-57; LIMA1, 0.88 versus 1.08, p = 8.1e-37), and left STAT6 null
(0.004, p = 0.67) and CDKN1A essentially unchanged (-0.13 versus -0.10,
p = 5.9e-20). Two estimates did move across a threshold: the TNFRSF1A null
result became nominally significant after correction (-0.013, p = 0.64 versus
-0.086, p = 0.028), and NCF2 reversed sign (-0.033, p = 0.16 versus +0.047,
p = 0.004), so no direction is interpreted for NCF2. For TNFRSF1B the global
test was not significant (p = 0.997), so no instrument was flagged and no
outlier-corrected estimate is available. Because cis-pQTL instruments come from a
single locus and remain partly correlated, both the MR-PRESSO global test and the
MR-Egger intercept test are anti-conservative here, and we therefore read them as
sensitivity checks rather than as evidence against a causal effect. Instrument
strength alone therefore does not explain the pattern; only CCM2 and LIMA1 carry
a coherent protein-level signal (Table~\ref{tab:mr_sensitivity}). Two of the nine
genes, \textit{TNF} and \textit{LTA}, lie within the MHC region (chr6:29--33 Mb),
where strong long-range LD violates the instrument-independence assumption of
every MR method used here; their estimates are therefore reported in
Table~\ref{tab:mr_sensitivity} for completeness and should be read as reference
values only, consistent with the fine-mapping caveat noted above.

\begin{table}[htbp]
\centering
\caption{Sensitivity analyses for the nine genes carrying cis-pQTL instruments.}
\label{tab:mr_sensitivity}
\footnotesize
\begin{tabular}{lllll}
\toprule
\multicolumn{5}{l}{\textbf{(A)} MR method comparison}\\
\midrule
  Gene & $n$ & IVW $b$ ($p$) & Weighted median $b$ ($p$) & Mode $b$ ($p$) \\
\midrule
  TNFRSF1A  & 25    & $-0.013$ (0.637)                & $+0.038$ (0.034)                 & $+0.030$ (0.075) \\
  NCF2      & 28    & $-0.033$ (0.156)                & $-0.105$ ($6.57\times10^{-6}$)   & $-0.115$ ($1.06\times10^{-4}$) \\
  STAT6     & 61    & $-0.007$ (0.427)                & $-0.012$ (0.082)                 & $-0.017$ ($7.37\times10^{-4}$) \\
  LIMA1     & 70    & $+0.882$ ($1.34\times10^{-26}$) & $+0.966$ ($7.03\times10^{-104}$) & $+0.909$ ($7.32\times10^{-66}$) \\
  TNFRSF1B  & 73    & $-0.051$ ($3.03\times10^{-7}$)  & $-0.042$ ($1.94\times10^{-3}$)   & $-0.040$ (0.096) \\
  CCM2      & 94    & $-0.515$ ($2.72\times10^{-17}$) & $-0.719$ ($7.26\times10^{-133}$) & $-0.734$ ($2.20\times10^{-44}$) \\
  CDKN1A    & 228   & $-0.131$ ($1.28\times10^{-18}$) & $-0.111$ ($1.86\times10^{-19}$)  & $-0.116$ ($1.41\times10^{-8}$) \\
  TNF$^{*}$ & 2288  & $-0.062$ ($8.14\times10^{-29}$) & $-0.014$ (0.085)                 & $+0.001$ (0.867) \\
  LTA$^{*}$ & 10213 & $-0.022$ ($5.41\times10^{-44}$) & $-0.012$ ($3.39\times10^{-22}$)  & $-0.001$ (0.907) \\
\bottomrule
\end{tabular}

\vspace{1.4em}

\begin{tabular}{lllll}
\toprule
\multicolumn{5}{l}{\textbf{(B)} Pleiotropy and heterogeneity diagnostics}\\
\midrule
  Gene & Egger intercept ($p$) & Cochran $Q$ $p$ & MR-PRESSO $p$ & MR-PRESSO distortion $p$ \\
\midrule
  TNFRSF1A  & $-0.0169$ ($5.13\times10^{-4}$)  & $2.28\times10^{-10}$  & $<2.0\times10^{-4}$ & 0.505 \\
  NCF2      & $+0.0409$ (0.034)                & $7.86\times10^{-11}$  & $2.0\times10^{-4}$ & $<2.0\times10^{-4}$ \\
  STAT6     & $+0.0205$ ($3.04\times10^{-12}$) & $1.74\times10^{-13}$  & $<2.0\times10^{-4}$ & 0.085 \\
  LIMA1     & $+0.0744$ (0.109)                & $5.28\times10^{-81}$  & $<2.0\times10^{-4}$ & $<2.0\times10^{-4}$ \\
  TNFRSF1B  & $-0.0138$ ($8.67\times10^{-5}$)  & 0.997                 & 0.997 & ---$^{\S}$ \\
  CCM2      & $+0.0129$ (0.549)                & $2.76\times10^{-107}$ & $<2.0\times10^{-4}$ & $<2.0\times10^{-4}$ \\
  CDKN1A    & $-0.0122$ ($1.76\times10^{-4}$)  & $8.78\times10^{-48}$  & $<2.0\times10^{-4}$ & $<2.0\times10^{-4}$ \\
  TNF$^{*}$ & $-0.0117$ ($<10^{-300}$)         & $<10^{-300}$          & skipped$^{\ddagger}$ & skipped$^{\ddagger}$ \\
  LTA$^{*}$ & $-0.0018$ ($7.41\times10^{-4}$)  & $<10^{-300}$          & skipped$^{\ddagger}$ & skipped$^{\ddagger}$ \\
\bottomrule
\end{tabular}

\vspace{1.4em}

\begin{tabular}{lll}
\toprule
\multicolumn{3}{l}{\textbf{(C)} Outlier handling}\\
\midrule
  Gene & Outliers & Outlier-corrected $b$ ($p$) \\
\midrule
  TNFRSF1A  & 5 & $-0.086$ (0.028) \\
  NCF2      & 3 & $+0.047$ (0.004) \\
  STAT6     & 5  & $+0.004$ (0.666) \\
  LIMA1     & 10 & $+1.076$ ($8.06\times10^{-37}$) \\
  TNFRSF1B  & none$^{\S}$ & ---$^{\S}$ \\
  CCM2      & 15 & $-0.762$ ($1.90\times10^{-57}$) \\
  CDKN1A    & 22 & $-0.102$ ($5.91\times10^{-20}$) \\
  TNF$^{*}$ & --- & skipped$^{\ddagger}$ \\
  LTA$^{*}$ & --- & skipped$^{\ddagger}$ \\
\bottomrule
\end{tabular}

\tblnote{\textbf{(A)} Effect estimates across MR methods, on the log-odds scale per 1 SD of genetically predicted protein level; genes are listed in ascending order of instrument count. \textbf{(B)} Pleiotropy and heterogeneity diagnostics: Cochran's $Q$ $p$-values test for heterogeneity across instruments; a non-significant MR-Egger intercept indicates no detectable directional pleiotropy; MR-PRESSO was run with \texttt{NbDistribution = 5000} so that the outlier test resolves $n/\mathrm{NbDistribution} \le 0.046$ for every gene tested, whereas at the package default of 1000 the same quantity is 0.061--0.228, above the 0.05 significance threshold, so the default cannot identify outliers at this instrument count. \textbf{(C)} Outlier handling: the outlier-corrected column reports the IVW estimate after removing the instruments flagged by the MR-PRESSO outlier test, and the Outliers column gives their number. Outlier-test $p$-values that fall at the resolution floor ($n/\mathrm{NbDistribution}$) are reported below as inequalities and are counted as outliers. For NCF2, removing 3 of the 28 instruments reverses the sign of the estimate (Table 3A versus 3C), so the direction of effect is not interpreted for that gene. $^{\S}$MR-PRESSO's outlier and distortion tests are performed only when the global test is significant, so for TNFRSF1B (p = 0.997) no instrument was tested for outlyingness and no outlier-corrected estimate exists. $^{*}$Located in the MHC region (chr6:29--33 Mb), where strong long-range LD violates the instrument-independence assumption of every method shown; estimates are reported for completeness and should be read as reference values only. $^{\ddagger}$Not computationally feasible: MR-PRESSO requires more than $n/0.046$ draws (about $5.0\times10^{4}$ for $n = 2288$ and $2.2\times10^{5}$ for $n = 10213$).}
\end{table}

Combined, none of the 4 UKB-PPP proteins and 1 of the 5 deCODE proteins with
cis-pQTL instruments showed pQTL--GWAS colocalization, compared with 57\% (43/76) at the eQTL
level (Fig.~\ref{fig:eqtl_pqtl_class}, \ref{fig:eqtl_pqtl_scatter}). Of the six non-MHC Tier 1 genes selected
for in-depth characterization (UTP11, BMP2, PNKD, LAMC1, TMBIM1, GPBAR1), only
BMP2 had a matching protein measurement (deCODE SomaScan), and it lacked
cis-pQTL instruments; the remaining five had no measurement in either panel.
The sole colocalizing protein, CCM2, is a distinct member of the 13-gene
non-MHC Tier 1 set, not one of the six core Tier 1 genes. The protein layer
therefore provides a systematic readout of transcript--protein divergence,
not direct orthogonal confirmation of the six core Tier 1 genes.

Steiger directionality testing across all nine genes with pQTL instruments
supported the instrument-to-protein direction. All 13,080 instrument-level
comparisons showed greater variance explained in plasma protein level than in
CRC risk, and the lead cis-pQTL variant was correctly oriented for every gene
(Steiger $P$ from $1.3\times10^{-4}$ for CCM2 to $P < 1\times10^{-300}$ for LTA;
Supplementary Table~S10). Directionality was therefore equally supported for the
single colocalizing gene and for the genes that did not colocalize, indicating
that the eQTL--pQTL discordance is not explained by reversed instrument
orientation.


\begin{figure}[htbp]
\centering
\begin{subfigure}{\textwidth}
\centering
\includegraphics[width=\linewidth]{Fig2a_coverage_funnel.pdf}
\caption{Protein-layer coverage funnel. Of the 62 protein-coding genes passing
Bonferroni-corrected SMR, 15 have plasma protein measurements in at least one of the two
proteomic panels (deCODE and UKB-PPP); adding the two TNF-pathway genes analyzed outside the
Bonferroni set (TNFRSF1A, TNFRSF1B) gives the 17-gene protein-layer analysis set. Seven of the
15 (nine of the 17) carry genome-wide significant cis-pQTL instruments, and one, CCM2, shows
pQTL--GWAS colocalization. Five of the six core Tier 1 candidates have no protein measurement in
either panel.}
\label{fig:coverage_funnel}
\end{subfigure}

\vspace{0.6em}

\begin{subfigure}{\textwidth}
\centering
\includegraphics[width=\linewidth]{Fig2b_eqtl_pqtl_classification.pdf}
\caption{Classification of the $n = 17$ genes with matched pQTL data by eQTL SMR and pQTL
Wald-ratio significance (two-sided $p < 0.05$): 4 concordant and significant (A), 1 discordant and
significant (B), and 12 without significant pQTL signal (C). Both axes use the single-variant
lead estimate. The one discordant gene is CDKN1A, and it is discordant by the Wald ratio only:
its multi-instrument IVW estimate is negative and therefore concordant with the eQTL layer
(Table~3). Two genes change category when the multi-instrument estimate is used instead---STAT6
loses nominal significance and TNFRSF1B gains it---so this panel should be read as a
single-variant view of the protein layer, not as the final classification.}
\label{fig:eqtl_pqtl_class}
\end{subfigure}

\vspace{0.6em}

\caption{The protein layer does not corroborate the genetic layer.}
\label{fig:protein_layer}
\end{figure}

\begin{figure}[tbp]
\ContinuedFloat
\centering
\setcounter{subfigure}{2}

\begin{subfigure}{\textwidth}
\centering
\includegraphics[width=\linewidth]{Fig2c_eqtl_pqtl_scatter.pdf}
\caption{Effect-size comparison between eQTL SMR and pQTL Mendelian randomization estimates for
the $n = 7$ genes with effect estimates in both layers. Six of seven genes showed concordant
effect direction; effect sizes were not significantly correlated (Spearman $\rho = -0.036$, $p = 0.96$).
Only CCM2 showed convergent evidence across both layers with coloc support.}
\label{fig:eqtl_pqtl_scatter}
\end{subfigure}
\end{figure}

\subsection{What can still be said: one convergent gene and a tissue-level example}

The one candidate that survives this coverage constraint is CCM2, which carries
convergent evidence across all three layers available in this study (Fig.~\ref{fig:ccm2_convergence}): eQTL
colocalization with a converged SuSiE credible set, a significant pQTL MR
association, and pQTL--GWAS colocalization. It is worth stating plainly that CCM2
belongs to the 13-gene non-MHC Tier 1 set but is not one of the six core Tier 1
genes selected for in-depth characterization: the gene with the strongest
cross-layer support is not the gene with the strongest genetic prioritization,
and we report it as such.

As a tissue-localization layer, we analyzed CRC single-cell RNA-seq data
($n = 50,214$ cells across 10 cell types) and spatial transcriptomics (Visium,
$n = 3,493$ spots). Of the 75 SMR genes, 61 (81\%) were detected in the
single-cell atlas; TMBIM1 showed fibroblast-enriched expression
(Supplementary Fig.~S14) and broad detection in CRC tissue (57\% of
3{,}493 Visium spots; Supplementary Fig.~S15). In contrast,
GPBAR1 was detected in only 1 of 3,493 Visium spots (0.03\%). Across the six core
Tier 1 candidates, therefore, five (UTP11, BMP2, PNKD, LAMC1 and TMBIM1) were
detectable in this section (11--57\% of spots; Supplementary Fig.~S15),
the mirror image of the protein layer, in which only BMP2 had a matching plasma
measurement and none of the six carried cis-pQTL instruments (Fig.~\ref{fig:coverage_funnel}). Partitioning
Visium spots into $k=6$ spatial microenvironment (ME) zones showed no zone-restricted expression: among the five core candidates with
detectable tissue expression (all but GPBAR1, which was detected in a single spot),
mean expression varied only 1.25--1.59-fold across the six zones, and UTP11, BMP2,
LAMC1 and TMBIM1 each reached their lowest mean expression in the stromal ME6
zone (PNKD: ME6 second lowest; Supplementary Fig.~S16;
Supplementary Table~S6). Pairwise spatial cross-correlation statistics for the
chr2:219Mb locus are provided in Supplementary Table~S7.

The tissue layer can localize a candidate gene within the tumor but cannot
validate causality. TCGA COAD+READ survival analysis~\cite{weinstein2013tcga} of the six non-MHC Tier 1 genes found no
significant associations (log-rank $p = 0.26$--$0.79$, all $p > 0.05$), and we report that negative result as-is rather than
presenting tissue expression as confirmation.

We evaluated directional consistency in FinnGen R13 (10,500 CRC cases), a
cohort with approximately 7.5-fold fewer cases than the discovery GWAS (10,500 vs 78,473) and therefore best
interpreted as a consistency check, not an adequately powered
replication sample. Among 8 coloc- or SMR-significant
genes with sufficient instrument availability (LAMC1, TMBIM1, PNKD, GPBAR1,
LIMA1, CASC8, UTP23, COX14; UTP11 and BMP2 were excluded due to allele
harmonization failure), only LIMA1 achieved nominal significance (IVW
$p = 0.0045$).
Five of the eight genes (LAMC1, TMBIM1, PNKD, GPBAR1, LIMA1) showed effect
directions concordant with the discovery analysis, whereas COX14, UTP23, and
CASC8 were direction-inconsistent.

To assess whether blood-based eQTL findings generalize to colon tissue, we
performed SMR using GTEx Colon Transverse cis-eQTL data. Of 25 tested genes
with colon cis-eQTLs at the standard threshold, 84\% showed concordant effect
direction with eQTLGen blood-based estimates (Spearman $\rho = 0.58$, $p = 0.0030$). All 10 genes
significant in both tissues showed identical direction. However, 5 of the 6 prioritized non-MHC Tier 1 genes lacked detectable colon cis-eQTLs, reflecting the power limitation of tissue-specific eQTL datasets (Fig.~\ref{fig:tissue_layer}). Forest plots for the FinnGen directional-consistency analysis are provided in Supplementary Fig.~S17.

\begin{figure}[htbp]
\centering
\begin{subfigure}{\textwidth}
\centering
\includegraphics[width=\linewidth]{Fig3a_ccm2_convergence.pdf}
\caption{Convergence at CCM2. Left: pQTL effect on CRC risk (Wald ratio $b$, 95\% CI) for the
five deCODE proteins with cis-pQTL instruments. Right: posterior probability of colocalization
between the cis-pQTL and CRC GWAS signals for the same five genes (dashed line, PPH4 = 0.8).
CCM2 is the only gene carrying both a genome-wide significant pQTL association and pQTL--GWAS
colocalization, and it is also the only gene in the study with eQTL--GWAS colocalization and
SuSiE fine-mapping support at the same locus.}
\label{fig:ccm2_convergence}
\end{subfigure}

\vspace{0.6em}

\begin{subfigure}{\textwidth}
\centering
\includegraphics[width=\linewidth]{Fig3b_gtex_colon_scatter.pdf}
\caption{Scatter of eQTLGen (blood) versus GTEx Colon Transverse SMR effect sizes for
$n = 25$ genes with colon cis-eQTL ($p < 5\times10^{-8}$); Spearman $\rho = 0.58$,
$p = 0.0030$, with 21 of 25 (84\%) direction-concordant.}
\label{fig:gtex_scatter}
\end{subfigure}

\vspace{0.6em}

\caption{GTEx Colon Transverse SMR sensitivity analysis.}
\label{fig:tissue_layer}
\end{figure}

\begin{figure}[tbp]
\ContinuedFloat
\centering
\setcounter{subfigure}{2}

\begin{subfigure}{\textwidth}
\centering
\includegraphics[width=\linewidth]{Fig3c_gtex_colon_forest.pdf}
\caption{Forest plot for the 10 genes with Bonferroni-significant SMR in both tissues.}
\label{fig:gtex_forest}
\end{subfigure}
\end{figure}

\subsection{From evidence to tractability}

To systematically assess target tractability, we classified the six non-MHC
Tier 1 genes according to the Illuminating the Druggable Genome (IDG)
framework (Pharos/DrugCentral; Tclin: approved drug exists; Tchem: chemical probe
or ligand available; Tbio: biological annotation only; Tdark: no functional
annotation). The six genes spanned a druggability gradient: LAMC1 (Tclin, ocriplasmin
approved), GPBAR1 (Tchem, GPCR with established ligand and chemical probe
infrastructure), BMP2 (Tchem, TGF-$\beta$ family with known receptor
inhibitors), UTP11 (Tchem, structure with ligand but zero antibody
tractability as a nucleolar protein), TMBIM1 (Tbio, no chemical probes or
known ligands---de novo discovery required despite highest genetic evidence),
and PNKD (Tbio, limited functional annotation with no drug discovery infrastructure; Supplementary Fig.~S18).

To assess potential safety risks of modulating these targets, we performed
PheWAS screening using Open Targets Platform genetic associations and summarized
each gene on two transparent dimensions rather than on a single composite score:
the number of distinct organ-system risk categories carrying at least one genetic
association, and the number of non-CRC high-risk phenotypes with genetic evidence
(Table~\ref{tab:tier1}; Supplementary Table~S8). BMP2 showed the narrowest risk
spectrum among the six non-MHC Tier 1 genes: its genetic associations were
confined to two categories (cardiovascular and neurological) and to two high-risk
phenotypes (atrial fibrillation; neurodegenerative disease). GPBAR1, by contrast,
had the fewest genetic associations of the six (5) but these included hepatic and
renal phenotypes consistent with its bile-acid physiology (8 high-risk
phenotypes). The remaining four genes spanned three to four categories each
(Fig.~\ref{fig:phewas}; Supplementary Fig.~S19). Because the composite
safety score in Supplementary Table~S8 and Supplementary Fig.~S19 is
dominated by the number of genetic associations, it ranks GPBAR1 numerically below
BMP2; we therefore report the two components separately in Table~\ref{tab:tier1}
and use those components, not the composite score, in the interpretation below.

Drug annotation using DrugBank and Open Targets identified 8 genes (13\% of
the 62 protein-coding set) with known pharmacological agents, including 4 with
approved drugs (Fig.~\ref{fig:drug}). LAMC1 was the highest-ranked drug candidate, with
the approved drug ocriplasmin (repositioning score = 0.756).
TNF had the largest drug portfolio (7 approved agents across 14 total entries
in DrugBank), matching its established role as a therapeutic target in
inflammatory disease. BMP2, a Tier 1 gene classified as Tchem by the IDG
framework, has an FDA-approved recombinant protein form (rhBMP-2/Infuse) used
therapeutically for bone regeneration.
Across the six core candidates, BMP2 is the only gene that combines strong
genetic support with an existing pharmacological precedent; for the others,
translation would require substantial additional drug-discovery work rather than
repositioning of an existing agent.

\begin{figure}[htbp]
\centering
\begin{subfigure}{\textwidth}
\centering
\includegraphics[width=\linewidth]{Fig4a_phewas_safety.pdf}
\caption{PheWAS screening of the six core non-MHC Tier 1 genes.
Heatmap of genetic disease associations across risk categories (Open Targets Platform).
Darker blue = more genetic evidence in that category. Among the six genes, the associations of
BMP2 are confined to the narrowest set of organ systems (cardiovascular and neurological) and to
the fewest high-risk phenotypes (2). The composite score $2.0\times$ risk-category diversity $+$
$1.5\times$ genetic evidence $+$ $0.5\times$ high-risk count is reported in Supplementary
Table~S8 and Supplementary Fig.~S19; it is dominated by the number of genetic
associations, so Table~\ref{tab:tier1} reports the two components separately and the main text
uses those.}
\label{fig:phewas}
\end{subfigure}

\vspace{0.8em}

\caption{Tractability and safety of the six core non-MHC Tier 1 candidates.}
\label{fig:tractability}
\end{figure}

\begin{figure}[tbp]
\ContinuedFloat
\centering
\setcounter{subfigure}{1}

\begin{subfigure}{\textwidth}
\centering
\includegraphics[width=\linewidth]{Fig4b_drug_repurposing.pdf}
\caption{Drug repurposing tier distribution.
Of $n = 62$ protein-coding SMR-significant genes, 8 (13\%) have known pharmacological agents
across DrugBank and Open Targets. Tier A: highest repurposing potential (coloc $+$ approved drug).}
\label{fig:drug}
\end{subfigure}
\end{figure}

%==========================================================================
\section{Discussion}
%==========================================================================

The evidence assembled here converges on a single point: for CRC, the genetic
layer of drug-target prioritization is resolved and close to saturated, whereas
the layer that would test those targets directly---plasma protein---is not. We
therefore frame these results not as another ranked target list but as a
quantification of where the evidence stops: which genes survive colocalization
and fine-mapping, which can be measured at the protein layer at all, and which
remain defensible once that coverage is made explicit.

\subsection*{Three explanations for the eQTL--pQTL gap}

Three non-mutually exclusive explanations could account for the systematic
discordance between eQTL and pQTL colocalization that we observe: limited
statistical power at the protein layer, post-transcriptional buffering, and
tissue mismatch between whole blood and the tissues of origin of circulating
proteins. While 57\% of SMR-significant genes (43/76) showed strong eQTL--GWAS
colocalization, none of the 5 UKB-PPP proteins and 1 of the 5 deCODE proteins
(CCM2, PPH4 = 0.9514) showed pQTL--GWAS colocalization. The protein-level
evidence should be read as validating the analytical pipeline and the existence
of systematic transcript--protein divergence, not as orthogonal confirmation of
the six non-MHC Tier 1 genes themselves: only BMP2 had a matching protein
measurement (deCODE SomaScan), and it lacked cis-pQTL instruments, precluding
protein-level validation. Their causal evidence therefore rests primarily on SMR,
coloc, and SuSiE fine-mapping, complemented by tissue-level (single-cell and
spatial) evidence, not protein-level support.

The first explanation is the limited statistical power of pQTL
studies. Although sample sizes are comparable ($n\approx35,000$ for both
deCODE and UKB-PPP vs. $n=31,684$ for eQTLGen), plasma proteomic
measurements have substantially greater technical variability than
transcriptomic measurements, reducing effective statistical power for
colocalization. This is compounded by the fact that only 15 of 62 protein-coding
genes had matching protein measurements, further constraining the pQTL
comparison. A second contributing factor is post-transcriptional regulatory
buffering: the well-established observation that mRNA levels explain only
30--40\% of protein abundance variance~\cite{vogel2012insights} suggests that
transcript-level genetic effects are systematically attenuated at the protein
level, which would predict lower colocalization rates at the pQTL layer. A
third factor is tissue specificity mismatch: eQTLGen captures whole-blood
eQTLs, while plasma pQTLs originate from multiple tissues (liver, immune cells,
intestinal epithelium), only partially overlapping with the causal CRC tissue.
Distinguishing among these explanations will require larger proteomic GWAS
with broader protein coverage, tissue-matched transcriptomic and proteomic
data from the same individuals, and formal power simulations for coloc at the
pQTL level. We therefore interpret the observed discordance as evidence that
\textit{current} pQTL datasets are insufficient for high-confidence drug
target colocalization at most loci, not as definitive proof of
biological divergence.

The apparent directional discordance at CDKN1A illustrates why the estimator
must be stated explicitly rather than why the biology is bidirectional. The eQTL
SMR estimate is protective ($b = -0.288$), and the single-variant lead pQTL Wald
ratio is nominally risk-increasing ($b = +0.130$, $p = 0.020$)---this is the one
discordant point in Fig.~\ref{fig:eqtl_pqtl_class}. The multi-instrument estimate at the same locus,
however, points the other way ($b = -0.131$, $p = 1.3\times10^{-18}$), as do the
weighted median and weighted mode estimates, and is therefore concordant with the
eQTL layer (Table~\ref{tab:mr_sensitivity}). The lead variant is rs1801270 (C31S),
a missense change that alters p21 protein stability, but the reversal of sign
between the single lead variant and the full set of 228 correlated cis variants
(Cochran's $Q$ $p = 8.8\times10^{-48}$) is a property of the instrument set, not
evidence that CDKN1A exerts opposite effects at transcript and protein level. We
therefore report CDKN1A as a locus whose direction depends on the estimator used,
and treat it as unresolved rather than as a biological finding. The same caution
applies in the opposite direction to TNFRSF1B, which is null by the Wald ratio but
significant by IVW: a single cis-pQTL lead variant is not a sufficient basis for
assigning direction, and multi-level evidence remains necessary before
drug-target prioritization.

\subsection*{Why the gap extends beyond CRC}

This gap is not specific to colorectal cancer. Every eQTL-based target list
inherits it, because the limiting factor is the availability of cis-pQTL
instruments rather than the strength of the genetic signal, and no current
plasma proteomic panel measures all protein-coding genes. The genetic layer
itself is close to saturated, and the incremental value of further SMR or coloc
screens therefore lies in the validation layers rather than in new gene lists.

Our gene set showed 44.9\% overlap with Chen et al.~2024, the
largest prior CRC genetic study (254K individuals, 289 resolvable gene symbols), and
17.4\% overlap with Hazelwood 2025 (52K meta-GWAS, 41 genes; Fig.~\ref{fig:competitor}).
Overlap with four additional low-tier studies (Hong et al.~\cite{hong2025integrative}, Wang et al.~\cite{wang2025discovoncol},
Xia and Wang~\cite{xia2025bmccancer}, Tian et al.~\cite{tian2025dbi}, based on
smaller FinnGen GWAS) was minimal
($\leq$3 genes total). The higher
overlap with Chen 2024 is expected given that both Chen 2024 and the present
study employ SMR from blood eQTL data, whereas Hazelwood 2025 relies on
multi-tissue TWAS with splicing QTLs---methodologically distinct approaches
that test different biological hypotheses and consequently identify partially non-overlapping gene sets~\cite{chen2024integrative, hazelwood2025multi}. The gene-level comparison underlying Fig.~\ref{fig:competitor} is provided in Supplementary Table~S9.
This is, to our knowledge, a systematic documentation of eQTL/pQTL
colocalization divergence in CRC susceptibility genetics, and it implies that
cis-pQTL coverage should be reported alongside any comparable target list.

\begin{figure}[htbp]
\centering
\begin{subfigure}{\textwidth}
\centering
\includegraphics[width=\linewidth]{Fig5a_competitor_overlap.pdf}
\caption{Gene overlap counts with six prior CRC genetic studies. Low-tier
studies (Hong, Wang, Wang BMC Cancer suppl., Tian) showed $\leq$3 gene overlap;
high-tier studies (Chen 2024, Hazelwood 2025) showed 44.9\% and 17.4\% overlap,
respectively.}
\label{fig:competitor_overlap}
\end{subfigure}

\vspace{0.6em}

\caption{The genetic layer is largely saturated: overlap with prior CRC target-discovery studies.}
\label{fig:competitor}
\end{figure}

\begin{figure}[tbp]
\ContinuedFloat
\centering
\setcounter{subfigure}{1}

\begin{subfigure}{\textwidth}
\centering
\includegraphics[width=\linewidth]{Fig5b_gene_overlap_venn.pdf}
\caption{Venn diagram of shared and unique gene symbols between this study (69 gene symbols),
Chen 2024 (289 resolvable gene symbols) and Hazelwood 2025 (41 genes). Circle areas are not
drawn to scale; every region is labelled with its exact symbol count and set sizes are given on
the right. Of this study's 69 genes, 34 (49\%) are not present in either prior gene list.}
\label{fig:competitor_upset}
\end{subfigure}
\end{figure}

Even among the overlapping genes, the analytical depth differs
substantially. The chr2:219Mb locus illustrates this distinction: both
Hazelwood 2025 (via coloc) and Chen 2024 (via conditional analysis) identified
signals at this region, whereas we examined the locus at variant level and asked
whether the three resident genes---PNKD, TMBIM1 and GPBAR1---carry separable cis
signals. They do not, with the reference panel available to us: the credible sets
returned for all three genes share the same lead variants (rs116204487 heads the
first credible set of each gene, and rs116515067 recurs in their second and third
sets), so the region does not resolve into three independent causal signals. This
matters for drug development: the three genes span a 222 kb region with distinct
druggability profiles (Tclin, Tbio, Tchem respectively), and the genetics cannot
currently tell the three candidate targets apart. UTP11 was not reported by any prior CRC study, and GPBAR1 was
absent from Chen 2024 but present in the Hazelwood 2025 druggable-genome set. CCM2, while present in both Chen 2024 and the Wang
BMC Cancer supplement, is the only gene across all studies to date to achieve
triple-convergent evidence (eQTL fine-mapping + pQTL MR + pQTL coloc), a level
of protein-level validation uniquely provided by our pipeline.

At the locus level, 23 of our 37 independent loci (62.2\%) contained genes
also reported by Chen 2024, while 14 loci (37.8\%) were fully independent in
our analysis. We interpret these data as indicating that the CRC susceptibility
transcriptome accessible through blood eQTL-based SMR has been substantially
mapped, and that the incremental contribution of future studies will come not
from discovering entirely new loci but from higher-resolution causal inference,
multi-omic integration (particularly the eQTL--pQTL gap we document), and
biologically contextualized validation. The distinct contribution of this
work is therefore the \textit{analytical depth} (variant-level LD-aware
fine-mapping used as a consistency check on colocalization, characterization of
the eQTL--pQTL colocalization discordance, spatial transcriptomics) and
\textit{target tractability assessment} (IDG druggability classification, PheWAS,
drug repurposing scoring) applied to both shared and unique loci.

\subsection*{What remains actionable}

Among our non-MHC Tier 1 genes, TMBIM1 had the most extensive multi-layer
support at the colocalization and spatial levels: coloc+SuSiE fine-mapping,
fibroblast-enriched expression in single-cell RNA-seq, and broad spatial
expression in CRC tissue (57\% of 3,493 Visium spots). The HEIDI test was
significant for TMBIM1 ($p_{\text{HEIDI}} = 1.62 \times 10^{-3}$), indicating
residual heterogeneity, so the multi-layer support is confined to
colocalization and spatial corroboration, not absence of pleiotropy.
TMBIM1 is not present
in the current deCODE or UKB-PPP pQTL panels, precluding protein-level
validation in our pipeline. TMBIM1 (transmembrane BAX inhibitor motif
containing 1) is a multi-pass transmembrane protein involved in
calcium-mediated apoptosis regulation. Currently there are no known chemical
probes or approved drugs targeting TMBIM1, and it is classified as
``Tbio'' (biological annotation only) by the Illuminating the Druggable
Genome (IDG) program, indicating that substantial drug discovery
infrastructure would need to be developed before it could enter preclinical
development. The fibroblast-enriched expression suggests a tumor
microenvironment targeting strategy that would be missed by bulk tissue
analyses alone, but its target tractability currently ranks below other Tier 1
candidates.

UTP11 is the only Tier 1 gene absent from all prior CRC studies. GPBAR1 was
absent from Chen 2024 but present in the Hazelwood 2025 druggable-genome set;
BMP2 was present in Chen 2024 but absent from Hazelwood 2025. LAMC1, PNKD, and
TMBIM1 were identified by both Hazelwood 2025 and Chen 2024. Independent
SuSiE fine-mapping resolved single-SNP credible sets at each of these loci; as
noted above, we read these as an LD-aware consistency check on colocalization
rather than as resolved causal variants. BMP2 (bone morphogenetic protein 2) is a member of the
TGF-$\beta$ superfamily with well-characterized roles in embryonic development
and bone morphogenesis, classified as Tchem by IDG with FDA-approved
recombinant protein therapy (Infuse) for bone regeneration. BMP2 showed the
largest per-SD CRC risk effect among all tested genes (OR = 1.38, 95\% CI:
1.27--1.50), exceeding the magnitude of established modifiable CRC risk
factors such as BMI and processed meat intake~\cite{papadimitriou2021umbrella}. BMP2 had
no HEIDI result because its probe was absent from the HEIDI output, a data gap,
not a passed heterogeneity test. In the context of
CRC, BMP2 signaling through the SMAD pathway intersects with the TGF-$\beta$
signaling axis frequently altered in colorectal tumorigenesis; BMP2 has been
reported to exert both tumor-suppressive and tumor-promoting effects depending
on cellular context and SMAD4 status. UTP11 (U3 small nucleolar
RNA-associated protein 11) is a ribosome biogenesis factor with no established
role in CRC biology and limited drug tractability (nucleolar protein,
inaccessible to antibody-based therapeutics), though its classification as
Tchem by IDG indicates at least structural ligand data.

GPBAR1 (Tchem), LAMC1 (Tclin), and PNKD (Tbio) span the full druggability
gradient among the chr2 and chr1 Tier 1 genes. GPBAR1 is the most druggable of
the three, as a G-protein-coupled bile acid receptor with established chemical
probe infrastructure (Tchem); however, its 25 disease phenotypes (5 with
genetic evidence; including cholesterol gallstone disease and NASH, mechanistically
linked to its bile acid physiology) constitute significant safety signals that
may reflect on-target physiology. LAMC1 is the only Tclin-classified gene
among our Tier 1 candidates, with the approved drug ocriplasmin targeting
laminin-mediated adhesion in vitreomacular adhesion; repositioning to CRC would
require substantial bridging pharmacology given that ocriplasmin is a
locally injected proteolytic enzyme and systemic laminin targeting carries
broad off-target toxicity risk. PNKD (paroxysmal nonkinesigenic dyskinesia) is
classified as Tbio with no known drugs or chemical probes and no detected
safety signals---representing the lowest tractability among Tier 1 genes
despite strong colocalization and SuSiE evidence. The HEIDI test was
significant for PNKD ($p_{\text{HEIDI}} = 7.22 \times 10^{-4}$), indicating
residual heterogeneity in the causal signal.

\subsection*{Limitations and methodological considerations}

Several limitations should be noted. First, our eQTL data derive from whole
blood, which may not fully capture colon-specific regulatory effects. Our GTEx
Colon Transverse sensitivity analysis showed 84\% direction concordance
overall (Spearman $\rho = 0.58$), providing moderate reassurance. However, 5 of the 6 prioritized non-MHC Tier 1 genes lacked detectable colon cis-eQTLs (at $p < 5\times10^{-8}$),
which could indicate either true blood-specific regulation or simply the power
limitation of GTEx's colon sample ($n\approx300$). Functional annotation of
colon-specific enhancers and chromatin state at these loci (e.g., from
ENCODE/Roadmap Epigenomics) could help resolve this ambiguity but was beyond
the scope of this study. Second, the FinnGen replication cohort is
substantially smaller than our discovery GWAS ($\sim$7.5$\times$ fewer cases),
comparable to the power gap between the present discovery GWAS (78K cases)
and the most recent prior study (Hazelwood 2025, 52K meta-GWAS; $\sim$1.5$\times$
power advantage). A formal power calculation (Supplementary Methods) indicates
that given the estimated effect sizes from discovery, the expected proportion
of nominally significant replications under a true-positive model is
$\approx$0.4, consistent with the observed 1/8 nominal significance. Third, we
did not perform wet-laboratory validation (RT-qPCR, IHC, or functional assays),
which would be needed to confirm the proposed causal mechanisms. Fourth, the
pQTL colocalization analysis was limited to genes with available matching
protein measurements (15 of 62 protein-coding genes); the PheWAS safety
screening used Open Targets Genetics associations, which reflect general
population GWAS signals, not drug-exposed adverse events, and should
not be interpreted as regulatory-grade safety assessment. Steiger testing supported the instrument-to-protein direction for every protein-layer instrument, but it cannot exclude reverse causation or horizontal pleiotropy; for LTA and TNF the aggregated test statistic was not numerically defined because thousands of correlated MHC-region variants were summed, so the direction call for those two genes rests on the per-instrument comparison. Fifth, we did not
stratify our analysis by CRC molecular subtype (CMS1--4, MSI/MSS status), and
the relevance of identified targets may differ across subtypes. Finally, while
87\% of protein-coding genes lacked known pharmacological agents, this reflects
both their recent GWAS-based identification and their limited tractability;
several top candidates (TMBIM1, PNKD) have no established chemical matter and
would require de novo drug discovery programs.

The tissue-localization layer carries additional limitations of its own. It rests on a
single Visium section ($n = 3{,}493$ spots; GSE285505, which contains further CRC
sections that were not analyzed here) and a single single-cell reference
($n = 50{,}214$ cells), so the observations above describe one tumor and were not tested
for replication in an independent section or cohort. These analyses were designed to ask
\emph{where} each candidate is expressed, not to estimate effect sizes: the
spatial--microenvironment associations are unadjusted for multiple comparisons. We therefore
treat the single-cell and spatial results as supporting
tissue-localization evidence rather than causal or functional evidence, and no independent
tissue validation (immunohistochemistry or in situ hybridization) was performed.

The observed genomic inflation factor ($\lambda = 1.600$) is elevated but
consistent with expectations for SMR/TWAS studies, where local LD structure
amplifies test statistic correlation compared to GWAS. In GWAS, the LD score
regression (LDSC) intercept provides a quantitative estimate of non-polygenic
inflation; however, LDSC is not directly applicable to SMR test statistics
because they are joint eQTL--GWAS tests, not single-trait GWAS summary
statistics. Elevated inflation is a recognized feature of SMR/TWAS test
statistics~\cite{wainberg2019, gusev2016integrative} and has been reported in
comparable CRC transcriptome-wide analyses~\cite{hazelwood2025multi}. We note
that the underlying GWAS (GCST90255675) has a large sample size
($n\approx185,000$) and is expected to show polygenic inflation, which would
propagate to SMR test statistics through shared LD architecture. The Bonferroni
correction we applied is conservative and controls the family-wise error rate
at the stated threshold, even in the presence of test correlation from local LD.
We did not derive an empirical null distribution for the SMR statistics (for
example by simulation under a matched LD structure), so the inflation factor is
reported as a descriptive observation rather than as a corrected test statistic.

%==========================================================================
\section{Conclusions}
%==========================================================================

This study shows that when CRC drug-target prioritization is driven by blood
eQTL-based SMR, the limiting step is protein-layer coverage rather than genetic
significance: of 15,582 genes tested, 75 reached SMR significance and 19 reached
Tier 1 status, yet only 15 of the 62 protein-coding genes were measured in either
plasma proteomic panel and only one colocalized at the protein layer. The
distinct contribution of this study
lies in analytical precision (variant-level LD-aware fine-mapping used as a
consistency check on colocalization, including the finding that the three genes
at the chr2:219Mb locus cannot be separated into independent cis signals with
current reference panels), multi-omic integration (a systematic documentation of
eQTL/pQTL colocalization discordance in CRC, with CCM2 as the sole
triple-convergent gene), and biological contextualization (spatial transcriptomics
characterizing tumor microenvironment zones and gene-level spatial
cross-correlations among Tier 1 genes). The eQTL/pQTL discordance suggests that current proteomic
GWAS datasets, while valuable, are insufficient for high-confidence
protein-level colocalization at most loci, and expanded proteomic resources
will be needed.

Once the coverage constraint is stated explicitly, the candidates that remain
defensible are CCM2 and BMP2: CCM2 carries the only triple-convergent evidence in
this study, and BMP2 combines the strongest genetic support with an existing
pharmacological precedent (a TGF-$\beta$ receptor-inhibitor class) and the
narrowest organ-system risk spectrum of the six core candidates (cardiovascular
and neurological phenotypes only; Table~\ref{tab:tier1}). We recommend that
cis-pQTL coverage be reported alongside any eQTL-derived target list, so that
readers can see where protein-level validation was possible and where it was not.

%==========================================================================
\bibliographystyle{unsrtnat}
\bibliography{manuscript/crc_smr_atlas}

%==========================================================================
%==========================================================================
\section*{Ethics Approval and Consent to Participate}
%==========================================================================

This study used only publicly available, de-identified summary-level data and
previously published de-identified single-cell and spatial transcriptomics data.
No new human participants or identifiable material were
collected, and institutional ethics approval was therefore not required. Consent
for publication is not applicable, as no individually identifiable data are
reported.

\section*{Funding}
%==========================================================================

This work was supported by the Jiangsu Provincial Health Commission Project
(Z2024049) and the Nanjing Medical University Science and Technology
Development Fund (NMUB20240108). The funders had no role in study design, data
collection and analysis, decision to publish, or preparation of the manuscript.

\section*{Author Contributions}

J.Z. is the first and corresponding author of this work.
%==========================================================================

J.Z. conceived and designed the study, performed all analyses, generated the
figures and tables, and wrote the manuscript. The author has read and approved
the final manuscript.

\section*{Competing Interests}
%==========================================================================

The author declares that he has no competing interests.

\section*{Acknowledgements}
%==========================================================================

The author thanks the investigators and participants of the resources used in
this study: eQTLGen, the Genotype-Tissue Expression (GTEx) Project, the UK
Biobank Pharma Proteomics Project, deCODE genetics, FinnGen, The Cancer Genome
Atlas, the NHGRI-EBI GWAS Catalog, and the Open Targets and Pharos teams.

\section*{Use of AI-Assisted Tools}
%==========================================================================

In accordance with the journal policy on AI-assisted technologies, the author
discloses that Claude (Anthropic) and Codex (OpenAI) were used for language
editing, code generation and debugging, and for a pre-submission simulated peer
review. All analytical decisions (statistical thresholds, parameter choices,
interpretation of results) were made by the author. All AI-generated code was
reviewed, tested and validated by the author, and every AI-suggested edit was
critically evaluated before acceptance. No generative AI tool was used to
fabricate or manipulate data, to create fraudulent references, or to circumvent
peer review. The author takes full responsibility for the accuracy, integrity
and originality of the work.
\section*{Data Availability}
%==========================================================================

The GWAS summary statistics for colorectal cancer (GCST90255675) are available
from the GWAS Catalog (\url{https://www.ebi.ac.uk/gwas/}). The eQTLGen
cis-eQTL data are available from \url{https://www.eqtlgen.org/}. The UKB-PPP
pQTL data are available from the UK Biobank Pharma Proteomics Project via
Synapse (project syn51364943) upon approved access. The deCODE pQTL data are
available from \url{https://download.decode.is/} upon request and approval.
The GTEx Colon Transverse eQTL data are available from the GTEx Portal
(\url{https://gtexportal.org/}). The FinnGen R13 summary statistics are
available from \url{https://www.finngen.fi/}. The 1000 Genomes Phase 3
reference panel is available from
\url{https://www.internationalgenome.org/}. TCGA COAD/READ expression data
were accessed via UCSC Xena (\url{https://xenabrowser.net/}). Single-cell
RNA-seq and Visium spatial transcriptomics data are available from GEO
(GSE285505). Chen et al.~2024 supplementary data were obtained from the
Nature Communications article supplement (\\allowbreak 41467\_2024\_47399\_MOESM4\_ESM.xlsx).

%==========================================================================
\section*{Code Availability}
%==========================================================================

All analysis scripts used in this study are openly available under the MIT
licence at \url{https://github.com/zjjwcwk2025/crc-smr-pqtl-atlas}. The repository
contains the complete analysis pipeline, the scripts that generate every figure
and table, and a README documenting the public data sources, software versions and
the order in which the documents are built. An archived copy of the scripts is
additionally available from the corresponding author.

\section*{Supplementary Information}
%==========================================================================

\begin{itemize}
\item Additional file 1 --- File name: Additional file 1; File format: PDF (.pdf);
Title: Supplementary figures and tables. Description: Supplementary Methods,
Supplementary Fig.~S1--S19 and Supplementary Table~S1--S10.
\item Additional file 2 --- File name: Additional file 2; File format: DOCX (.docx);
Title: Supplementary Tables S1--S10 (editable). Description: Editable versions of
Supplementary Tables S1--S10 for production.
\item Additional file 3 --- File name: Additional file 3; File format: PDF (.pdf);
Title: STROBE-MR checklist. Description: Completed STROBE-MR checklist for this
Mendelian randomization analysis, item by item, with the location of each item in
the manuscript.
\end{itemize}

\end{document}
